\section{A final audit table}
\label{sec:second-incompleteness-05}

\begin{longtable}{@{}p{0.24\textwidth}p{0.32\textwidth}p{0.34\textwidth}@{}}
\caption{A final audit table}\\
\toprule
Claim & Required hypotheses & Common invalid strengthening \\
\midrule
\endfirsthead
\caption[]{A final audit table (continued)}\\
\toprule
Claim & Required hypotheses & Common invalid strengthening \\
\midrule
\endhead
\midrule\multicolumn{3}{r}{\emph{continued on next page}}\\
\endfoot
\endlastfoot
First incompleteness, G\"odel sentence & \(T\) computably enumerable,
\(T\supseteq\mathsf Q\); consistency for unprovability; omega-consistency or
\(\Sigma^0_1\)-soundness for independence of \(G_T\), that is, for
underivability of both \(G_T\) and \(\neg G_T\) & Every consistent theory
has the displayed G\"odel sentence independent under no effectiveness or
arithmetic-strength assumption. \\
\addlinespace
First incompleteness, Rosser & \(T\) consistent, computably enumerable, and
\(T\supseteq\mathsf Q\) & Every consistent theory whatever, with no
computable-enumerability or \(\mathsf Q\)-extension hypothesis, has an
independent sentence. \\
\addlinespace
Second incompleteness & \(T\supseteq\mathsf Q\) consistent and computably
enumerable; the previously fixed provability predicate satisfies all three
derivability conditions &
No theory can prove a consistency statement about itself or any weaker theory.
\\
\addlinespace
L\"ob's theorem & Diagonal lemma and all three derivability conditions & Any
informal predicate called ``provable'' satisfies L\"ob's theorem. \\
\addlinespace
Tarski undefinability & One arithmetic formula purporting to define truth for
all arithmetic sentences in \(\mathcal N\) & No restricted satisfaction or
partial truth predicate can be defined. \\
\bottomrule
\end{longtable}

